% chapter: 参考文献
\chapter{参考文献}\label{章：参考文献}

本書に関係する書籍と,本文に対応する個別ウェブ資料をまとめた.
ウェブ資料は説明・公式の照合や補足学習の参照先である.
本書と資料で型の名前・初項の番号・\bookterm{normalize}{正規化}が異なる場合があるため,
対応箇所を手掛かりに,条件と記法を照合して読んでほしい.

\begingroup
\raggedright
\section{書籍}
\begin{enumerate}[label={[\arabic*]}]
  \item \begin{minipage}[t]{\linewidth}\raggedright
  松坂和夫,\,『新装版 数学読本1』,\,岩波書店,\,2019年.
  ISBN 978-4-00-029877-3.
  \par\noindent\url{https://www.hanmoto.com/bd/isbn/9784000298773}\par
  \end{minipage}\par\smallskip

  \item \begin{minipage}[t]{\linewidth}\raggedright
  松坂和夫,\,『新装版 数学読本2』,\,岩波書店,\,2019年.
  ISBN 978-4-00-029878-0.
  \par\noindent\url{https://www.books.or.jp/book-details/9784000298780}\par
  \end{minipage}\par\smallskip

  \item \begin{minipage}[t]{\linewidth}\raggedright
  松坂和夫,\,『新装版 数学読本3』,\,岩波書店,\,2019年.
  ISBN 978-4-00-029879-7.
  \par\noindent\url{https://www.hanmoto.com/bd/isbn/9784000298797}\par
  \end{minipage}\par\smallskip

  \item \begin{minipage}[t]{\linewidth}\raggedright
  松坂和夫,\,『新装版 数学読本4』,\,岩波書店,\,2019年.
  ISBN 978-4-00-029880-3.
  \par\noindent\url{https://www.hanmoto.com/bd/isbn/9784000298803}\par
  \end{minipage}\par\smallskip

  \item \begin{minipage}[t]{\linewidth}\raggedright
  松坂和夫,\,『新装版 数学読本5』,\,岩波書店,\,2019年.
  ISBN 978-4-00-029881-0.
  \par\noindent\url{https://www.hanmoto.com/bd/isbn/9784000298810}\par
  \end{minipage}\par\smallskip

  \item \begin{minipage}[t]{\linewidth}\raggedright
  松坂和夫,\,『新装版 数学読本6』,\,岩波書店,\,2019年.
  ISBN 978-4-00-029882-7.
  \par\noindent\url{https://www.hanmoto.com/bd/isbn/9784000298827}\par
  \end{minipage}\par\smallskip

  \item \begin{minipage}[t]{\linewidth}\raggedright
  宮腰忠,\,『高校数学+α：基礎と論理の物語』,\,共立出版,\,2004年.
  ISBN 978-4-320-01768-9.
  \par\noindent\url{https://www.kyoritsu-pub.co.jp/book/b10010328.html}\par
  \end{minipage}\par\smallskip

\end{enumerate}

\section{本編の漸化式：おいしい数学}
\begin{enumerate}[resume,label={[\arabic*]}]
  \item \begin{minipage}[t]{\linewidth}\raggedright
  おいしい数学,\,「\bookterm{recurrence}{漸化式}の全パターン紹介」.
  \par{\small 対応：第3章の全体像.\par}
  \par\noindent\url{https://hiraocafe.com/note/zenkashiki.html}\par
  \noindent（2026年10月9日閲覧）.
  \end{minipage}\par\smallskip

  \item \begin{minipage}[t]{\linewidth}\raggedright
  おいしい数学,\,「2-1型（等差型）の\bookterm{recurrence}{漸化式}」.
  \par{\small 対応：基本4パターン：等差型.\par}
  \par\noindent\url{https://hiraocafe.com/note/zenkashiki2-1.html}\par
  \noindent（2026年10月9日閲覧）.
  \end{minipage}\par\smallskip

  \item \begin{minipage}[t]{\linewidth}\raggedright
  おいしい数学,\,「2-2型（等比型）の\bookterm{recurrence}{漸化式}」.
  \par{\small 対応：基本4パターン：等比型.\par}
  \par\noindent\url{https://hiraocafe.com/note/zenkashiki2-2.html}\par
  \noindent（2026年10月9日閲覧）.
  \end{minipage}\par\smallskip

  \item \begin{minipage}[t]{\linewidth}\raggedright
  おいしい数学,\,「2-3型（\bookterm{difference}{階差}型）の\bookterm{recurrence}{漸化式}」.
  \par{\small 対応：基本4パターン：\bookterm{difference}{階差}型.\par}
  \par\noindent\url{https://hiraocafe.com/note/zenkashiki2-3.html}\par
  \noindent（2026年10月9日閲覧）.
  \end{minipage}\par\smallskip

  \item \begin{minipage}[t]{\linewidth}\raggedright
  おいしい数学,\,「2-8型（\bookterm{ratio}{階比}型）の\bookterm{recurrence}{漸化式}」.
  \par{\small 対応：基本4パターン：\bookterm{ratio}{階比}型／\bookterm{auxiliary}{補助数列}の\bookterm{normalize}{正規化}.\par}
  \par\noindent\url{https://hiraocafe.com/note/zenkashiki2-8.html}\par
  \noindent（2026年10月9日閲覧）.
  \end{minipage}\par\smallskip

  \item \begin{minipage}[t]{\linewidth}\raggedright
  おいしい数学,\,「2-4型（\bookterm{characteristic}{特性方程式}型）の\bookterm{recurrence}{漸化式}」.
  \par{\small 対応：\bookterm{characteristic}{特性方程式}型／\bookterm{auxiliary}{補助数列}の設計.\par}
  \par\noindent\url{https://hiraocafe.com/note/zenkashiki2-4.html}\par
  \noindent（2026年10月9日閲覧）.
  \end{minipage}\par\smallskip

  \item \begin{minipage}[t]{\linewidth}\raggedright
  おいしい数学,\,「2-5型（指数型）の\bookterm{recurrence}{漸化式}」.
  \par{\small 対応：\bookterm{difference}{階差}等比複合型：指数的な付加項.\par}
  \par\noindent\url{https://hiraocafe.com/note/zenkashiki2-5.html}\par
  \noindent（2026年10月9日閲覧）.
  \end{minipage}\par\smallskip

  \item \begin{minipage}[t]{\linewidth}\raggedright
  おいしい数学,\,「2-7型（関数スライド型）の\bookterm{recurrence}{漸化式}」.
  \par{\small 対応：\bookterm{difference}{階差}等比複合型：多項式の付加項／\bookterm{auxiliary}{補助数列}の設計.\par}
  \par\noindent\url{https://hiraocafe.com/note/zenkashiki2-7.html}\par
  \noindent（2026年10月9日閲覧）.
  \end{minipage}\par\smallskip

  \item \begin{minipage}[t]{\linewidth}\raggedright
  おいしい数学,\,「2-6型（逆数型）の\bookterm{recurrence}{漸化式}」.
  \par{\small 対応：分数型：逆数型.\par}
  \par\noindent\url{https://hiraocafe.com/note/zenkashiki2-6.html}\par
  \noindent（2026年10月9日閲覧）.
  \end{minipage}\par\smallskip

  \item \begin{minipage}[t]{\linewidth}\raggedright
  おいしい数学,\,「2-9型（\bookterm{logarithm}{対数}型）の\bookterm{recurrence}{漸化式}」.
  \par{\small 対応：\bookterm{logarithm}{対数}型.\par}
  \par\noindent\url{https://hiraocafe.com/note/zenkashiki2-9.html}\par
  \noindent（2026年10月9日閲覧）.
  \end{minipage}\par\smallskip

  \item \begin{minipage}[t]{\linewidth}\raggedright
  おいしい数学,\,「2-10型（1次分数型）の\bookterm{recurrence}{漸化式}」.
  \par{\small 対応：分数型：1次分数型.\par}
  \par\noindent\url{https://hiraocafe.com/note/zenkashiki2-10.html}\par
  \noindent（2026年10月9日閲覧）.
  \end{minipage}\par\smallskip

  \item \begin{minipage}[t]{\linewidth}\raggedright
  おいしい数学,\,「3-1型の\bookterm{recurrence}{漸化式}（隣接3項間\bookterm{recurrence}{漸化式}の基本）」.
  \par{\small 対応：隣接3項間\bookterm{recurrence}{漸化式}型.\par}
  \par\noindent\url{https://hiraocafe.com/note/zenkashiki3-1.html}\par
  \noindent（2026年10月9日閲覧）.
  \end{minipage}\par\smallskip

  \item \begin{minipage}[t]{\linewidth}\raggedright
  おいしい数学,\,「3-2型の\bookterm{recurrence}{漸化式}（隣接3項間\bookterm{recurrence}{漸化式}の応用）」.
  \par{\small 対応：隣接3項間\bookterm{recurrence}{漸化式}型の拡張.\par}
  \par\noindent\url{https://hiraocafe.com/note/zenkashiki3-2.html}\par
  \noindent（2026年10月9日閲覧）.
  \end{minipage}\par\smallskip

  \item \begin{minipage}[t]{\linewidth}\raggedright
  おいしい数学,\,「和を含んだ\bookterm{recurrence}{漸化式}」.
  \par{\small 対応：総和型.\par}
  \par\noindent\url{https://hiraocafe.com/note/zenkashikisn.html}\par
  \noindent（2026年10月9日閲覧）.
  \end{minipage}\par\smallskip

  \item \begin{minipage}[t]{\linewidth}\raggedright
  おいしい数学,\,「連立\bookterm{recurrence}{漸化式}」.
  \par{\small 対応：連立型.\par}
  \par\noindent\url{https://hiraocafe.com/note/zenkashikirenritsu.html}\par
  \noindent（2026年10月9日閲覧）.
  \end{minipage}\par\smallskip

\end{enumerate}

\section{本編と付録：高校数学の美しい物語}
\begin{enumerate}[resume,label={[\arabic*]}]
  \item \begin{minipage}[t]{\linewidth}\raggedright
  学びTimes「高校数学の美しい物語」,\,「\bookterm{recurrence}{漸化式}の\bookterm{characteristic}{特性方程式}の意味とうまくいく理由」.
  \par{\small 対応：\bookterm{characteristic}{特性方程式}型／\bookterm{auxiliary}{補助数列}の設計／隣接多項間.\par}
  \par\noindent\url{https://manabitimes.jp/math/1299}\par
  \noindent（2026年10月9日閲覧）.
  \end{minipage}\par\smallskip

  \item \begin{minipage}[t]{\linewidth}\raggedright
  学びTimes「高校数学の美しい物語」,\,「三項間\bookterm{recurrence}{漸化式}の3通りの解き方」.
  \par{\small 対応：隣接3項間\bookterm{recurrence}{漸化式}型／付録の\bookterm{matrix}{行列}による見直し.\par}
  \par\noindent\url{https://manabitimes.jp/math/697}\par
  \noindent（2026年10月9日閲覧）.
  \end{minipage}\par\smallskip

  \item \begin{minipage}[t]{\linewidth}\raggedright
  学びTimes「高校数学の美しい物語」,\,「一次分数型の\bookterm{recurrence}{漸化式}の解法と例題」.
  \par{\small 対応：分数型：1次分数型.\par}
  \par\noindent\url{https://manabitimes.jp/math/1331}\par
  \noindent（2026年10月9日閲覧）.
  \end{minipage}\par\smallskip

  \item \begin{minipage}[t]{\linewidth}\raggedright
  学びTimes「高校数学の美しい物語」,\,「\bookterm{matrix}{行列}の\bookterm{eigenvalue}{固有値}・\bookterm{eigenvector}{固有ベクトル}の定義と具体的な計算方法」.
  \par{\small 対応：付録：\bookterm{matrix}{行列}の基本、図5.2.\par}
  \par\noindent\url{https://manabitimes.jp/math/1008}\par
  \noindent（2026年10月9日閲覧）.
  \end{minipage}\par\smallskip

  \item \begin{minipage}[t]{\linewidth}\raggedright
  学びTimes「高校数学の美しい物語」,\,「\bookterm{matrix}{行列}の対角化の意味と具体的な計算方法」.
  \par{\small 対応：付録：\bookterm{characteristic}{特性方程式}を\bookterm{matrix}{行列}で見直す.\par}
  \par\noindent\url{https://manabitimes.jp/math/1133}\par
  \noindent（2026年10月9日閲覧）.
  \end{minipage}\par\smallskip

  \item \begin{minipage}[t]{\linewidth}\raggedright
  学びTimes「高校数学の美しい物語」,\,「\bookterm{chebyshev}{チェビシェフ多項式}」.
  \par{\small 対応：虚数解型隣接3項間／付録：\bookterm{chebyshev}{チェビシェフ多項式}.\par}
  \par\noindent\url{https://manabitimes.jp/math/729}\par
  \noindent（2026年10月9日閲覧）.
  \end{minipage}\par\smallskip

  \item \begin{minipage}[t]{\linewidth}\raggedright
  学びTimes「高校数学の美しい物語」,\,「\bookterm{sequence}{数列}の\bookterm{generating}{母関数}の意味とその応用例」.
  \par{\small 対応：付録：\bookterm{sequence}{数列}全体と\bookterm{generating}{母関数}.\par}
  \par\noindent\url{https://manabitimes.jp/math/658}\par
  \noindent（2026年10月9日閲覧）.
  \end{minipage}\par\smallskip

  \item \begin{minipage}[t]{\linewidth}\raggedright
  学びTimes「高校数学の美しい物語」,\,「カタラン数の意味と\bookterm{recurrence}{漸化式}」.
  \par{\small 対応：付録：\bookterm{convolution}{畳み込み}とカタラン数.\par}
  \par\noindent\url{https://manabitimes.jp/math/657}\par
  \noindent（2026年10月9日閲覧）.
  \end{minipage}\par\smallskip

  \item \begin{minipage}[t]{\linewidth}\raggedright
  学びTimes「高校数学の美しい物語」,\,「オイラー法をわかりやすく解説」.
  \par{\small 対応：付録：一歩の変化と微積分.\par}
  \par\noindent\url{https://manabitimes.jp/math/2536}\par
  \noindent（2026年10月9日閲覧）.
  \end{minipage}\par\smallskip

  \item \begin{minipage}[t]{\linewidth}\raggedright
  学びTimes「高校数学の美しい物語」,\,「ガンマ関数（階乗の一般化）の定義と性質」.
  \par{\small 対応：付録：階乗からガンマ関数へ.\par}
  \par\noindent\url{https://manabitimes.jp/math/960}\par
  \noindent（2026年10月9日閲覧）.
  \end{minipage}\par\smallskip

\end{enumerate}

\section{公式と条件：NIST DLMF}
\begin{enumerate}[resume,label={[\arabic*]}]
  \item \begin{minipage}[t]{\linewidth}\raggedright
  NIST Digital Library of Mathematical Functions,\,「§3.8 Nonlinear Equations」.
  \par{\small 対応：\bookterm{newton}{ニュートン法}型／付録の近似計算.\par}
  \par\noindent\url{https://dlmf.nist.gov/3.8}\par
  \noindent（2026年10月9日閲覧）.
  \end{minipage}\par\smallskip

  \item \begin{minipage}[t]{\linewidth}\raggedright
  NIST Digital Library of Mathematical Functions,\,「§3.3 Interpolation」.
  \par{\small 対応：付録の\bookterm{interpolation}{ラグランジュ補間}コラム.\par}
  \par\noindent\url{https://dlmf.nist.gov/3.3}\par
  \noindent（2026年10月9日閲覧）.
  \end{minipage}\par\smallskip

  \item \begin{minipage}[t]{\linewidth}\raggedright
  NIST Digital Library of Mathematical Functions,\,「§18.9 Recurrence Relations and Derivatives」.
  \par{\small 対応：付録のチェビシェフ・\bookterm{legendre}{ルジャンドル多項式}.\par}
  \par\noindent\url{https://dlmf.nist.gov/18.9}\par
  \noindent（2026年10月9日閲覧）.
  \end{minipage}\par\smallskip

  \item \begin{minipage}[t]{\linewidth}\raggedright
  NIST Digital Library of Mathematical Functions,\,「§18.3 Definitions」.
  \par{\small 対応：付録の\bookterm{legendre}{ルジャンドル多項式}の直交性.\par}
  \par\noindent\url{https://dlmf.nist.gov/18.3}\par
  \noindent（2026年10月9日閲覧）.
  \end{minipage}\par\smallskip

  \item \begin{minipage}[t]{\linewidth}\raggedright
  NIST Digital Library of Mathematical Functions,\,「§18.12 Generating Functions」.
  \par{\small 対応：付録の\bookterm{chebyshev}{チェビシェフ多項式}の\bookterm{generating}{母関数}.\par}
  \par\noindent\url{https://dlmf.nist.gov/18.12}\par
  \noindent（2026年10月9日閲覧）.
  \end{minipage}\par\smallskip

  \item \begin{minipage}[t]{\linewidth}\raggedright
  NIST Digital Library of Mathematical Functions,\,「§5.2 Definitions」.
  \par{\small 対応：付録のガンマ関数の定義.\par}
  \par\noindent\url{https://dlmf.nist.gov/5.2}\par
  \noindent（2026年10月9日閲覧）.
  \end{minipage}\par\smallskip

  \item \begin{minipage}[t]{\linewidth}\raggedright
  NIST Digital Library of Mathematical Functions,\,「§5.5 Functional Relations」.
  \par{\small 対応：付録のガンマ関数の更新.\par}
  \par\noindent\url{https://dlmf.nist.gov/5.5}\par
  \noindent（2026年10月9日閲覧）.
  \end{minipage}\par\smallskip

\end{enumerate}

\section{整数の更新と数え上げ}
\begin{enumerate}[resume,label={[\arabic*]}]
  \item \begin{minipage}[t]{\linewidth}\raggedright
  Keith Conrad,\,\textit{Pell's Equation I}.
  \par{\small 対応：\bookterm{pell}{ペル方程式}の具体例・整数解を作る更新.\par}
  \par\noindent\url{https://kconrad.math.uconn.edu/blurbs/ugradnumthy/pelleqn1.pdf}\par
  \noindent（2026年10月10日閲覧）.
  \end{minipage}\par\smallskip

  \item \begin{minipage}[t]{\linewidth}\raggedright
  Keith Conrad,\,\textit{Pell's Equation II}.
  \par{\small 対応：\bookterm{pell}{ペル方程式}と\bookterm{continued}{連分数}の近似分数の関係.\par}
  \par\noindent\url{https://kconrad.math.uconn.edu/blurbs/ugradnumthy/pelleqn2.pdf}\par
  \noindent（2026年10月10日閲覧）.
  \end{minipage}\par\smallskip

  \item \begin{minipage}[t]{\linewidth}\raggedright
  NIST Digital Library of Mathematical Functions,
  \,$\S$26.3 \textit{Lattice Paths: Binomial Coefficients}.
  \par{\small 対応：\bookterm{double}{二重数列}の格子経路,二項係数とその加法公式.\par}
  \par\noindent\url{https://dlmf.nist.gov/26.3}\par
  \noindent（2026年10月10日閲覧）.
  \end{minipage}\par\smallskip

\end{enumerate}
\endgroup
